%! AP Statistics — Unit 2, Lesson 2.5 — Homework (Answer Key)
\documentclass[10pt]{article}
\usepackage{apstats-article}
\usepackage{apstats-key}

\begin{document}

\pageheader{Unit 2, Lesson 2.5}{Homework --- Answer Key}
\namedateperiod

\vspace{0.16in}

\noindent\textbf{Directions:} Complete all problems. Write all interpretations in complete
sentences with context. Show technology results where requested.

\vspace{0.14in}

% Problem 1 ──────────────────────────────────────────────────────────────────
\noindent\textbf{1.} For each correlation coefficient, describe the direction and strength
of the linear association and give a real-world example that could produce that $r$ value.

\vspace{8pt}
\renewcommand{\arraystretch}{1.7}
\begin{tabularx}{\linewidth}{@{} C{1.4cm} X X @{}}
\rowcolor{sky}
\textbf{$r$} & \textbf{Direction \& Strength} & \textbf{Real-world example} \\
\midrule
$0.95$  & \ans{Strong positive linear} & \ans{Hours studied vs.\ exam score} \\[2pt]
$-0.72$ & \ans{Strong negative linear} & \ans{Engine size vs.\ fuel efficiency} \\[2pt]
$0.08$  & \ans{Very weak / no linear association} & \ans{Shoe size vs.\ GPA} \\[2pt]
$-0.99$ & \ans{Very strong negative linear} & \ans{Temperature vs.\ hot-drink sales} \\[2pt]
$0.45$  & \ans{Moderate positive linear} & \ans{Sleep hours vs.\ mood score} \\
\end{tabularx}

\vspace{0.2in}

% Problem 2 ──────────────────────────────────────────────────────────────────
\noindent\textbf{2.} Match each scatterplot description (A -- E) to the $r$ value that best
fits. Each value is used exactly once.

\vspace{8pt}
\textbf{$r$ values:} $\quad -0.93 \qquad -0.42 \qquad 0.00 \qquad 0.68 \qquad 0.99$

\vspace{8pt}
\renewcommand{\arraystretch}{1.5}
\begin{tabularx}{\linewidth}{@{} l X c @{}}
\rowcolor{sky}
\textbf{Plot} & \textbf{Description} & \textbf{$r$} \\
\midrule
A & Points scattered randomly with no pattern & \ans{$0.00$} \\
B & Points tightly clustered on a line sloping upward to the right & \ans{$0.99$} \\
C & Points loosely follow a line sloping upward & \ans{$0.68$} \\
D & Points tightly follow a line sloping steeply downward & \ans{$-0.93$} \\
E & Points loosely follow a downward pattern with several outliers & \ans{$-0.42$} \\
\end{tabularx}

\vspace{0.2in}

% Problem 3 ──────────────────────────────────────────────────────────────────
\noindent\textbf{3.} A study finds that the number of fire trucks sent to a fire ($x$) is
strongly positively correlated with the amount of fire damage ($y$), with $r = 0.92$.

\begin{enumerate}[label=\textbf{(\alph*)}, leftmargin=2em, itemsep=0.16in]
  \item Write a complete interpretation of $r = 0.92$ in context.\\[8pt]
        \ansline{There is a strong positive linear association between the number of fire
        trucks dispatched and the amount of fire damage (in \$1000s) for these fires
        ($r = 0.92$). Fires that receive more trucks tend to have greater damage.}\\[3pt]
        \ansline{}
  \item A city council member says this shows that sending more fire trucks causes more
        damage. Explain why this conclusion is not justified. Identify a likely lurking
        variable.\\[8pt]
        \ansline{This conclusion is not justified because correlation does not imply causation.
        The size of the fire is a lurking variable: larger fires cause more damage \emph{and}
        require more trucks. The trucks are responding to the fire, not causing the damage.}\\[3pt]
        \ansline{}
\end{enumerate}

\vspace{0.2in}

% Problem 4 ──────────────────────────────────────────────────────────────────
\noindent\textbf{4.} The table below shows data for 8 athletes: daily protein intake in grams
($x$) and muscle gain in pounds over 8 weeks ($y$).

\vspace{8pt}
\renewcommand{\arraystretch}{1.4}
\begin{tabularx}{\linewidth}{@{} l *{8}{C{1cm}} @{}}
\rowcolor{sky}
\textbf{Protein (g)} & 80 & 100 & 110 & 130 & 140 & 150 & 170 & 180 \\
\textbf{Gain (lb)}   & 1.2 & 1.8 & 2.1 & 2.9 & 3.3 & 3.5 & 4.0 & 4.4 \\
\end{tabularx}

\begin{enumerate}[label=\textbf{(\alph*)}, leftmargin=2em, itemsep=0.16in]
  \item Use technology to compute $r$. \quad $r = \ans{0.999}$
  \item Write a complete interpretation of $r$ in context.\\[8pt]
        \ansline{There is a very strong positive linear association between daily protein
        intake (grams) and muscle gain (pounds) over 8 weeks for these 8 athletes
        ($r \approx 1.00$). Athletes who consume more protein tend to gain more muscle.}\\[3pt]
        \ansline{}
  \item Does the correlation prove that eating more protein causes muscle gain?
        Explain.\\[8pt]
        \ansline{No. This is an observational study, so we cannot establish causation from
        $r$ alone. Athletes who eat more protein may also train harder or sleep better
        --- these lurking variables could explain the muscle gains.}\\[3pt]
        \ansline{}
\end{enumerate}

\vspace{0.18in}

\begin{extensionbox}
\small\textbf{Extension --- Optional}

Look up ``spurious correlations'' online (e.g., the website by Tyler Vigen). Choose one example
with a high $r$ value. Write a paragraph explaining: (a) what the two variables are,
(b) what the $r$ value is, and (c) why this correlation almost certainly does not reflect causation.
\end{extensionbox}

\begin{teachernote}
Extension answers will vary. Classic examples include per-capita cheese consumption vs.\
deaths by bedsheet tangling ($r \approx 0.95$) and Nicolas Cage films released vs.\
pool drownings. Credit any answer that correctly identifies the two variables, states $r$,
and articulates why the association is coincidental or driven by a lurking variable (often
time --- both series rise or fall over the same years).
\end{teachernote}

\vspace{0.14in}

\begin{tcolorbox}[colback=greenbg, colframe=greenacc, boxrule=0.8pt, arc=2mm,
    left=4mm, right=4mm, top=3mm, bottom=3mm]
\small\textbf{Preview --- Lesson 2.6}\\[4pt]
In Lesson 2.6 we will use $r$ to build the \textbf{least-squares regression line} --- a line
that uses the correlation to make predictions. The slope of that line depends directly on $r$.

\vspace{4pt}
\textbf{Think ahead:} If $r = 0$, what would the slope of the regression line be? What
would you predict for $y$ regardless of $x$?\\[6pt]
\ansline{If $r = 0$, the slope equals $0$. The regression line would be horizontal, and
the predicted value of $y$ would be $\bar{y}$ (the mean of $y$) for every $x$.}\\[2pt]
\ansline{}
\end{tcolorbox}

\end{document}
